\section{Comparison with naive approach}
\label{app:comparison}

\begin{algorithm}[t]
\caption{Naive private gradient oracle}
\label{alg:oracle-standard}
\begin{algorithmic}[1]
    \Statex \textbf{Input:} dataset $\calZ$, constant $B$, noise $\sigma$
    \Statex \textbf{Initialize:} Partition $\calZ$ into $KT$ subsets $Z_t^k$ of size $B$. The data size is $|\calZ|=M=BKT$.
    \State Upon receiving round index $k,t$ and parameter $\vecw_t^k$:
    \State Sample $\vecu_1,\ldots,\vecu_B\sim\mathrm{Uniform}(\bfS)$ i.i.d.
    \State $\vecg_t^k \gets \frac{1}{B}\sum_{i=1}^B \frac{d}{2\delta}(f(\vecw_t^k+\delta\vecu_i,z_i) - f(\vecw_t^k-\delta\vecu_i,z_i))\vecu_i$ where $z_i$ is the $i$-th data in $Z_t^k$.
    \State Return $\tilde\vecg_t^k \gets \vecg_t^k + \xi_t^k$ where $\xi_t^k \sim \calN(0,\sigma^2I)$ i.i.d.
\end{algorithmic}
\end{algorithm}

To concrete the discussion from Section \ref{sec:algorithm}, we'll delve into the analysis of the naive private algorithm built on O2NC. Specifically, this algorithm replaces the gradient oracle (as seen in line 5 of Algorithm \ref{alg:PONC}) with a direct zeroth-order estimator, as defined in Algorithm \ref{alg:oracle-standard}. As a remark, $\vecg_t^k$ is an unbiased estimator and $\E\|\vecg_t^k\|^2 \lesssim L^2(\frac{d}{B}+1)$, according to \citep{kornowski2023algorithm} Remark 6. Moreover, following the same argument as the sensitivity bound in Lemma \ref{lem:grad-estimator}, the sensitivity of $\vecg_t^k$ is bounded by $O(dL/B)$. Consequently, if we set $\sigma=dL/B\rho$, each $\tilde\vecg_t^k$ achieves $(\alpha,\alpha\rho^2/2)$-RDP. 

Although the algorithm appears more straightforward, its convergence analysis is less favorable. Following the proof idea of Theorem \ref{thm:convergence}, we first recall the result in Corollary \ref{cor:O2NC-smoothing}:
\begin{equation*}
    \E\|\nabla F(\overline \vecw)\|_{2\delta} 
    \le \frac{F^*+2L\delta}{DKT} + \sum_{k=1}^K\frac{\E\regret_T(\hat\vecu^k)}{DKT} + \sum_{k=1}^K\sum_{t=1}^T \frac{ \E\langle \nabla\hat F_\delta(\vecw_t^k)-\tilde\vecg_t^k, \Delta_t^k-\hat\vecu^k\rangle }{DKT}.
\end{equation*}
Recall that $\E[\regret_T(\hat\vecu^k)] \le D\sqrt{\sum_{t=1}^T \E\|\tilde\vecg_t^k\|^2}$. Therefore, we can bound the regret of OSD by
\begin{align*}
    \E[\regret_T(\hat\vecu^k)]
    &\le \textstyle D\sqrt{\sum_{t=1}^T \E\|\vecg_t^k\|^2 + \E\|\xi_t^k\|^2} \\
    &\lesssim \textstyle D\sqrt{\sum_{t=1}^T L^2(\frac{d}{B}+1) + \frac{d^3L^2}{B^2\rho^2}} \\
    &\lesssim \textstyle DL\sqrt{T} \left( \frac{\sqrt{d}}{\sqrt{B}} + 1 + \frac{d^{3/2}}{B\rho} \right).
\end{align*}
Next, we bound the sum of inner product. It's important to note that $\E\langle \nabla\hat F_\delta(\vecw_t^k)-\tilde\vecg_t^k, \Delta_t^k \rangle = 0$ because $\Delta_t^k$ is independent of $Z_t^k$. Therefore, using the martingale concentration bound, we have
\begin{align*}
    \textstyle
    \sum_{t=1}^T \E\langle \nabla\hat F_\delta(\vecw_t^k)-\tilde\vecg_t^k, \Delta_t^k-\hat\vecu^k \rangle 
    &= \textstyle
    \E\langle \sum_{t=1}^T \nabla\hat F_\delta(\vecw_t^k)-\tilde\vecg_t^k, -\hat\vecu^k \rangle \\
    &\le \textstyle 
    D\sqrt{\sum_{t=1}^T \E\|\nabla\hat F_\delta(\vecw_t^k)-\vecg_t^k\|^2 + \E\|\xi_t^k\|^2} \\
    &\lesssim \textstyle 
    DL\sqrt{T} \left( \frac{\sqrt{d}}{\sqrt{B}} + 1 + \frac{d^{3/2}}{B\rho} \right).
\end{align*}
% The second line follows from the martingale bound. 
Upon substituting back into Corollary \ref{cor:O2NC-smoothing}, we have
\begin{align*}
    \E\|\nabla F(\overline \vecw)\|_{2\delta} 
    &\lesssim
    \frac{F^*+L\delta}{DKT} + \frac{L}{\sqrt{T}}\left(\frac{\sqrt{d}}{\sqrt{B}}+1+\frac{d^{3/2}}{B\rho}\right)
    \intertext{Upon replacing $D=\delta/T$ and $M=KBT$, we have}
    &\lesssim 
    \frac{(F^*+L\delta)BT}{\delta M} + \frac{\sqrt{d}L}{\sqrt{BT}} + \frac{L}{\sqrt{T}} + \frac{d^{3/2}L}{B\sqrt{T}\rho}
    \intertext{To achieve the optimal non-private rate $O(d\delta^{-1}\epsilon^{-3})$, we need to set $B=1$. Therefore, upon setting $T=\min\{(\frac{\sqrt{d}L\delta M}{F^*+L\delta})^{2/3}, (\frac{d^{3/2}L\delta M}{(F^*+L\delta)\rho})^{2/3}\}$, we have}
    &\lesssim \left(\frac{(F^* + L\delta)dL^2}{\delta M}\right)^{1/3} + \left(\frac{(F^*+L\delta)d^{3/2}L^2}{\rho\delta M}\right)^{1/3}.
\end{align*}
Equivalently, $\E\|\nabla F(\overline\vecw)\|_{2\delta}\le \epsilon$ if $M=\Omega\left( d(\frac{F^*L^2}{\delta\epsilon^3}+\frac{L^3}{\epsilon^3}) + d^{3/2}(\frac{F^*L^2}{\rho\delta\epsilon^3} + \frac{L^3}{\rho\epsilon^3}) \right)$, whose dominating term is $\Omega(d^{3/2}\rho^{-1}\delta^{-1}\epsilon^{-3})$. In contrast, our carefully designed algorithm only requires $\tilde\Omega(d\delta^{-1}\epsilon^{-3}+d^{3/2}\rho^{-1}\delta^{-1}\epsilon^{-2})$ samples, as proved in Theorem \ref{thm:convergence}. Notably, the naive algorithm suffers an additional cost of $\epsilon^{-1}$ for privacy. 
% Moreover, note that the private term dominates the naive bound, suggesting that the ``privacy for free'' property when $\rho\ge \sqrt{d}\epsilon$ no longer holds.